\section{Weighted consequences}\label{ap:weighted}

The density bound gives uniform bounds on harmonic sums and tails over
progression-free sets, and permits weights larger than the reciprocal.
It also implies progressions in every subset of the primes with positive
relative upper density.  We derive these consequences from dyadic
summation and the resulting counting estimates.

Fix $k\ge3$. In this section write $c=c_k$, $C=C_k$, and
$\varepsilon=\varepsilon_k$ for constants such that
\begin{equation}\label{ap:weighted:density}
 r_k(N)\le CN\exp\bigl(-c(\log N)^\varepsilon\bigr)
 \qquad(N\ge2).
\end{equation}
We may decrease $\varepsilon$ and enlarge $C$ so that
$0<\varepsilon<1$. All bounds below are uniform over sets containing
no nonconstant $k$-term arithmetic progression.

\subsection{Dyadic summation}

\begin{proposition}[Dyadic weighted summation]\label{ap:weighted:dyadic}
For $w:\mathbb N\to[0,\infty)$, set
\[
 S_k(w)=\sum_{m=0}^{\infty}r_k(2^m)
                  \max_{2^m\le n<2^{m+1}}w(n).
\]
If $A\subseteq\mathbb N$ contains no nonconstant $k$-term
progression, then
\[
                       \sum_{a\in A}w(a)\le S_k(w).
\]
Consequently, whenever $S_k(w)<\infty$, divergence of the weighted
sum forces a nonconstant $k$-term progression.
\end{proposition}

\begin{proof}
Each interval $[2^m,2^{m+1})\cap\mathbb Z$ has length $2^m$,
so translation invariance gives at most $r_k(2^m)$ elements of $A$
there. Bound each weight by the maximum on that interval and sum.
All terms are nonnegative, so passage to infinite sums is valid.
\end{proof}

\begin{corollary}[Uniform harmonic bound]\label{ap:weighted:harmonic}
There is $H_k<\infty$ such that every $k$-term-progression-free
$A\subseteq\mathbb N$ satisfies
\[
              \sum_{a\in A}\frac1a\le H_k.
\]
One may take $H_k=\sum_{m\ge0}2^{-m}r_k(2^m)$.
\end{corollary}

\begin{proof}
Apply Proposition~\ref{ap:weighted:dyadic} with $w(n)=1/n$.
By Equation~\eqref{ap:weighted:density}, its terms for $m\ge1$ are at most
$C\exp(-c(m\log2)^\varepsilon)$, whose sum converges.
\end{proof}

\subsection{Weights larger than the reciprocal}

The following weights are defined by their formulas on a sufficiently
large tail and set to zero before that tail. Changing finitely many
values does not affect a divergence criterion.

\begin{corollary}[Stretched-exponential weights]\label{ap:weighted:stretched}
Every $k$-term-progression-free set has a uniformly bounded sum for
each of the weights
\begin{align*}
 w_a(n)&=\frac{\exp(a(\log n)^\varepsilon)}{n},
                  &&0<a<c,\\
 w_{*,\eta}(n)&=
 \frac{\exp(c(\log n)^\varepsilon)}
 {n\log n\,(\log\log n)^{1+\eta}},
                  &&\eta>0.
\end{align*}
The bound may depend on $k$ and the displayed parameters. Divergence
for either weight therefore forces a nonconstant $k$-term progression.
\end{corollary}

\begin{proof}
Both weights are eventually decreasing. Indeed, with $y=\log n$,
the derivatives of their logarithms with respect to $y$ are
\[
 -1+a\varepsilon y^{\varepsilon-1},\qquad
 -1+c\varepsilon y^{\varepsilon-1}
       -y^{-1}-(1+\eta)(y\log y)^{-1},
\]
respectively. Since $\varepsilon<1$, both are negative for large $y$.
On every sufficiently late dyadic interval the maximum is therefore
its left-endpoint value. For $w_a$, the corresponding summand in
Proposition~\ref{ap:weighted:dyadic} is at most
\[
                    C e^{-(c-a)(m\log2)^\varepsilon},
\]
which is summable. For $w_{*,\eta}$ the exponential factors cancel
at that same endpoint, leaving
\[
          \frac{C}{m\log2\,[\log(m\log2)]^{1+\eta}},
\]
also summable. The omitted initial intervals contribute a finite amount.
\end{proof}

\subsection{Logarithmic weights and dense subsets of the primes}

The exponent $\varepsilon$ in the preceding two weights depends on the fixed
length $k$. Applying Theorem~\ref{ap:main} separately for each fixed
length gives the following consequence with a single logarithmic weight.

\begin{corollary}[Logarithmic weights and every length]
\label{ap:weighted:logarithmic}
For every fixed $B\ge0$, if $A\subseteq\mathbb N$ satisfies
\[
            \sum_{a\in A}\frac{(\log(2+a))^B}{a}=\infty,
\]
then $A$ contains nonconstant arithmetic progressions of every finite
length.
\end{corollary}

\begin{proof}
Fix $k\ge3$. On the $m$th dyadic interval the weight is at most
$C_B(1+m)^B2^{-m}$. Hence the terms in
Proposition~\ref{ap:weighted:dyadic} are bounded, apart from finitely
many, by
\[
           C_{k,B}(1+m)^B e^{-c_k(m\log2)^{\varepsilon_k}}.
\]
This sequence is summable because $m^{\varepsilon_k}/\log m\to\infty$.
Thus a set with the stated divergence cannot be $k$-term-progression-free.
Lengths one and two follow from infinitude.
\end{proof}

The same comparison shows that, for each fixed $B>0$, there is a
constant $C_{k,B}$ with
\[
               r_k(N)\le C_{k,B}N(\log N)^{-B}\qquad(N\ge3).
\]

\paragraph{Dense subsets of the primes.}
Theorem~\ref{ap:main} also recovers the established dense-primes theorem
of Green and Tao~\cite[Theorem~1.2]{GreenTaoPrimes2008}.
If $A$ is a subset of the primes with
$\limsup_{N\to\infty}|A\cap[N]|/\pi(N)>0$, where $\pi(N)$ counts
the primes at most $N$, then $A$ contains infinitely many nonconstant
$k$-term progressions for every fixed $k\ge3$.
Indeed, the preceding estimate with $B=2$ gives
$r_k(N)=o(N/\log N)$. Positive relative upper density and the prime
number theorem give $c>0$ and an unbounded sequence of $N$ with
$|A\cap[N]|\ge cN/\log N$. For each fixed $M$, deleting $A\cap[M]$
leaves at least $(c/2)N/\log N$ elements along all sufficiently large
members of this sequence. This exceeds $r_k(N)$, so there is a
progression entirely above $M$. Since $M$ is arbitrary, infinitely
many occur. Lengths one and two follow from infinitude.

\subsection{Uniform harmonic tails}

\begin{corollary}[Harmonic tails]\label{ap:weighted:tail}
There is $C'_k>0$ such that every $k$-term-progression-free set
$A\subseteq\mathbb N$ satisfies, for all sufficiently large real $x$,
\[
 \sum_{\substack{a\in A\\a>x}}\frac1a
 \le C'_k(\log x)^{1-\varepsilon_k}
                e^{-c_k(\log x)^{\varepsilon_k}}.
\]
In particular, for every $0<a<c_k$, this tail is at most
$C_{k,a}e^{-a(\log x)^{\varepsilon_k}}$.
\end{corollary}

\begin{proof}
Let $A(t)=|A\cap[1,t]|$. The function
$t e^{-c(\log t)^\varepsilon}$ is eventually increasing, since
$\varepsilon<1$. Applying Equation~\eqref{ap:weighted:density} at
$\lfloor t\rfloor$ therefore gives
$A(t)\le Ct e^{-c(\log t)^\varepsilon}$ for all sufficiently large
real $t$. Partial summation and $A(t)/t\to0$ yield
\[
 \sum_{\substack{a\in A\\a>x}}\frac1a
 \le C\int_x^\infty e^{-c(\log t)^\varepsilon}\frac{dt}{t}
 =\frac{C c^{-1/\varepsilon}}{\varepsilon}
       \int_{c(\log x)^\varepsilon}^\infty
                  v^{1/\varepsilon-1}e^{-v}\,dv.
\]
Put $b=1/\varepsilon-1>0$. If $v_0\ge2b$, then
$(v/v_0)^b\le e^{(v-v_0)/2}$ for $v\ge v_0$, by
$\log(1+t)\le t$. Thus
\[
       \int_{v_0}^\infty v^b e^{-v}\,dv
                         \le2v_0^b e^{-v_0}.
\]
Substitution proves the first claim, with constant $2C/(c\varepsilon)$
on a sufficiently large tail. Finally the polynomial factor in
$\log x$ is absorbed by
$e^{(c-a)(\log x)^\varepsilon}$, proving the second claim.
\end{proof}
